\section{Reference Manual}
\label{app:reference}

\subsection{Public Contract}

\texttt{Jsonifiable} packages an encoder, decoder, and one checked theorem:

\begin{lstlisting}[language=Coq]
Class Jsonifiable (A : Type) := {
  to_JSON     : A -> JSON;
  from_JSON   : JSON -> Result A string;
  canonical_jsonification : forall (a : A), from_JSON (to_JSON a) = res a
}.
\end{lstlisting}

The theorem states encoder-to-decoder preservation.  It does not state that
every JSON value accepted by the decoder is the unique canonical representation
for the type.

\subsection{Supported Declarations}

The deriver targets computational declarations in \texttt{Type} or
\texttt{Set}.  It supports ordinary enumerations, algebraic datatypes, records,
parameterized types whose parameters have \texttt{Jsonifiable} instances, and
structural recursion over the target.  The current implementation also supports
a conservative set of nested recursive occurrences through standard containers
such as lists, options, and products.

Unsupported declarations fail without registering a partial instance.  Current
non-targets or limitations include targets in \texttt{Prop}, dependent
constructor fields, indexed families, proof fields, function-valued fields,
sort-valued fields, coinductive declarations, arbitrary nested recursion, and
mutual recursion.

\subsection{Encoding Conventions}

\begin{center}
  \small
  \begin{tabular}{p{0.40\linewidth}p{0.58\linewidth}}
    \toprule
    Rocq shape                     & JSON shape                                                                                \\
    \midrule
    Nullary constructor            & JSON string containing the constructor name.                                              \\
    Constructor with fields        & Singleton JSON object keyed by constructor name.                                          \\
    Record                         & JSON object keyed by projection names.                                                    \\
    Named constructor argument     & Field key is the binder name.                                                             \\
    Anonymous constructor argument & Field key is derived from the argument type head and disambiguated with numeric suffixes. \\
    Type parameter                 & Instance argument requiring a \texttt{Jsonifiable} dictionary.                            \\
    \bottomrule
  \end{tabular}
\end{center}

Primitive and library instances can choose representation-specific encodings.
For example, \texttt{nat} is encoded as a JSON number rather than as the unary
constructor structure of Rocq natural numbers.

\subsection{Generated Artifacts}

For a declaration \texttt{T}, a successful command
\texttt{Elpi derive.jsonifiable T} adds definitions for the encoder, decoder,
round-trip proof, and registered instance.  The visible instance name follows
the generated-name convention used in the examples, such as
\texttt{UserRole\_Jsonifiable} or \texttt{BinTree\_Jsonifiable}.  Users should
normally rely on typeclass resolution and the public operations shown in the
\texttt{Jsonifiable} class rather than depending on internal helper names.

\subsection{Interpreting Failure Categories}

The benchmark artifact records failures with categories intended for triage.
Some categories are semantic non-targets for a generic JSON serializer:
propositions are erased at extraction, arbitrary functions do not have a
canonical finite representation, sorts are type-level objects, and coinductive
values may be infinite.  Other categories are implementation limitations:
dependent fields, indexed families, proof-bearing computational wrappers,
arbitrary nested recursion, and mutual recursion require additional contracts or
proof-generation techniques.
